\noindent\textbf{GAN-based audio generation.} GAN-based models consist of two networks: a generator that produces audio waveforms and a discriminator that classifies whether the generated audio is real or synthetic. The key to these methods is designing sophisticated discriminators to capture audio features from different perspectives, thereby guiding the generator to produce more fine-grained and realistic audio through adversarial learning. 
MelGAN~\citep{melgan} employs a Multi-Scale Discriminator (MSD) that extracts features from waveforms at various downsampled resolutions. HiFi-GAN~\citep{hifigan} extends this approach by incorporating a Multi-Period Discriminator (MPD) that learns complementary structures by analyzing different periodic patterns. Univnet~\citep{mrd} utilizes a Multi-Resolution Discriminator (MRD) that operates in the time-frequency domain across multiple spectral resolutions. BigVGAN~\citep{bigvgan} introduces the Snake activation function to provide periodic inductive bias and incorporates low-pass filters for anti-aliasing. Unlike previous works that rely on transposed convolutions for upsampling, Vocos~\citep{vocos} models spectral coefficients using ConvNeXt model~\citep{convnext} and realizes upsampling to waveforms through the Inverse Short-Time Fourier Transform (ISTFT). Our network structure follows Vocos by using ConvNeXt to operate on spectral coefficients, but we extend it to multiple time-frequency resolutions for a more powerful backbone.

\noindent\textbf{Diffusion/Flow Matching-based audio generation.}
Diffusion-based methods learn transition trajectories from noise distribution to data distribution, excelling at stable training and robust generation. 
DiffWave~\citep{diffwave} and WaveGrad~\citep{Wavegrad} first introduce diffusion models to audio synthesis and get GAN-competitive results. PriorGrad~\citep{priorgrad} uses an adaptive prior derived from data statistics to structure the noise, thereby speeding up the diffusion sampling process. 
PeriodWave~\citep{periodwave} utilizes a multi-period Flow Matching estimator to learn different periodic features of the waveform.
Multi-Band Diffusion (MBD)~\citep{mbd} divides the audio into different frequency bands and processes each band with a distinct diffusion model. RFWave~\citep{rfwave} also adopt this multi-band idea, employing Flow Matching to model each band but processing all bands in parallel with a shared model. These multi-band processing approaches employ frequency equalization to reduce the energy mismatch between the Gaussian distribution and audio data distribution in different frequency bands. While this technique is effective for high-frequency modeling, the equalization process relies on normalization statistics calculated on the dataset~\citep{mbd} or accumulated during training~\citep{rfwave}, which poses a potential generalization risk when applied to distributions that differ significantly from those used to compute the original statistics. 
In this work, we improve Flow Matching for audio modeling by reformulating the objective to endpoint estimation and emphasizing loss in perceptually salient regions, and further combine a GAN fine-tuning stage for fast and finer-grained generation. Note that PriorGrad and RFWave also use energy-based loss scaling~\citep{rfwave}, but they operates at the per-frame level, whereas ours is based on both time and frequency dimensions, which is validated to outperform the per-frame version in our reformulated setting (Table~\ref{tab:ablation_flow}).

\noindent\textbf{Accelerating Flow Matching models in audio generation.} A typical way to accelerate diffusion model inference is through distillation~\citep{progressive_distillation,guided_distillation,consistency_model}, using multi-step outputs to guide fewer-step outputs. WaveFM~\citep{wavefm} introduces a tailored method of consistency distillation~\citep{consistency_model} and can generate audio in 1 step without significant quality degradations. 
Another approach is to integrate GAN objectives with adversarial learning using real data~\citep{fastdiff2, periodwave_turbo, revisiting_distillation}. 
% A recent work in audio generation is PeriodWave-Turbo~\citep{periodwave_turbo}, which introduces GAN fine-tuning to Flow-Matching pre-trained models. Our Flow2GAN shares a similar idea, but we build upon our enhanced Flow Matching specialized for audio modeling, and our 2-step generator achieves competitive results compared to their 4-step model, while offering much faster inference.
A recent work in audio generation is PeriodWave-Turbo~\citep{periodwave_turbo}, which introduces GAN fine-tuning to Flow-Matching pre-trained models. Our Flow2GAN shares a similar idea, but we build upon our enhanced Flow Matching specialized for audio modeling, providing stronger pretrained weights to initialize the GAN generators. This results in significantly superior one-step generator after GAN fine-tuning compared to building upon standard Flow Matching (Table~\ref{tab:ablation_flow}). Notably, PeriodWave-Turbo~\citep{periodwave_turbo} does not explore one-step results. 
Shortcut models~\citep{shortcut} condition on both noise level and step size, enabling better few-step generation than consistency models~\citep{consistency_model}. We also investigate this strategy for few-step audio generation and demonstrate that our approach surpasses it.



